\subsection{Roster Notation} 
Previously, we listed a set
\[
A=\{1,2,3\}
\]
This set is given in \emph{roster notation}.
\begin{definition}[Roster Notation]
A set is listed in \term{roster notation} by listing all of its elements, enclosed in \term{set brackets}.
\end{definition}
\begin{Note}Set brackets are curly. Here's how to write them.
\begin{cpic}{}{2}
\end{cpic}
\end{Note}
For larger sets, it's not always practical (or possible!) to list all of the elements. In this case, we use \makeblank to shortcut how they're written.
\begin{Note}
The \makeblank in which we write the elements of the set still doesn't matter, it is important to see the pattern.
\end{Note}
\begin{Example}
The set $\{1,2,3,\dots,10\}$ consists of \makeblank[1in] numbers from \makeblanksmall to \makeblanksmall. If we wrote out the entire set, it would be:
\[
\makeblank[3in]
\]
\end{Example}
\begin{Example}
The set $\{2,4,6,\dots,100\}$ consists of \makeblank[1in] numbers from \makeblanksmall to \makeblanksmall. (We won't write out the entire set here.)
\end{Example}
We can also use three dots at the end to mean the set ``goes on forever.''
\begin{Example}
The set $\{3,6,9,12,\dots\}$ contains \emph{all} the numbers that we get counting by \makeblank[1in].
\end{Example}
It should be \emph{easy} to make sense of the pattern.
%\begin{Example}
%If we wrote $\{3,6,9,\dots,100\}$, the pattern doesn't make sense:
%\begin{itemize}
%\item We are counting by \makeblank[1in]
%\item But if we continue that pattern, we would count \makeblanksmall followed by \makeblanksmall and miss $100$ entirely!
%\end{itemize}
%If we rewrite this set as
%\[
%\makeblank[3in]
%\]
%we have a set that makes sense.
%\end{Example}
\begin{Example}
If we write $\{1,3,7,\dots\}$ it isn't obvious what the pattern is.
\begin{itemize}
\item The pattern might be: add \makeblanksmall, add \makeblanksmall, so next we would add \makeblanksmall to get \makeblanksmall as our next number.
\item The pattern might also be: \makeblank[1in] the number add add \makeblanksmall, so the next number would be \makeblank.
\item Maybe you can spot other possible patterns, too!
\end{itemize}
In a case like this, we might be better off to use \emph{set-builder notation}, which we'll meet later.
\end{Example}